\section{Full electric sources and physical preparation}\label{sec:full-electric-preparation}
We now test two further mechanisms in the same finite slab and
state: regular-point distributional sources for the full electric
operator, and observables passed through the strict half-slab
integral. Both give positive actual pairings. Neither realizes Q.
No continuum assumption, physical mass gap, or RH is used.

\subsection{Sobolev powers of full-electric point sources}
There are $L=648$ spatial links and $V=216$ spatial vertices.
The compact link manifold has dimension $D=3L=1944$.
Choose an irreducible connection q: after a spanning-tree gauge
choice, two noncommuting loop holonomies have common centralizer
$\{\pm I\}$. The principal gauge quotient near q is smooth of
dimension
\begin{equation}\label{eq:quotient-dimension}
d=3L-3V=1296.
\end{equation}
The scalar and traceless covariant sectors are locally bundles of
ranks one and three. Both are actual orthogonal reducing subspaces
of $E=E_\nu$; the center acts trivially on their fibers. Their
electric restrictions $E_j$ are locally elliptic with the same
scalar positive quadratic principal symbol $p_2(q,\eta)$.

Write $E_je_{jn}=\lambda_{jn}e_{jn}$ in actual orthonormal bases,
$\omega_{jn}=\sqrt{1+\lambda_{jn}}$, and choose nonzero fiber
vectors $\xi_j$ at q independently of arithmetic. For any real s
define the combined spectral source
\begin{equation}\label{eq:full-point-source}
J_sf=\sum_{j=0}^1\sum_n\omega_{jn}^{-s}b_{jn}
 \widehat f(\epsilon_j\omega_{jn})e_{jn},\qquad
b_{jn}=\ip{e_{jn}(q)}{\xi_j}_{\rm fib},\quad
\epsilon_0=1,\ \epsilon_1=-1.
\end{equation}
Point preparation fixes a boundary configuration; its electric
Sobolev power controls regularity. These definitions precede any
arithmetic matching. The two signs use existing color sectors;
the generator is $\sqrt{1+E_0}\oplus(-\sqrt{1+E_1})$, an electric
wave action rather than physical time.

The local spectral law at this regular point is
\begin{equation}\label{eq:point-weyl}
N_j(R):=\sum_{\omega_{jn}\le R}|b_{jn}|^2
=C_jR^d+o(R^d),\qquad C_j>0.
\end{equation}
In a local orthonormal fiber frame with quotient density
$a_\nu(q)\dd q$,
\begin{equation}
C_j=\frac{\norm{\xi_j}_{\rm fib}^2}{(2\pi)^d a_\nu(q)}
\int_{p_2(q,\eta)<1}\dd\eta.
\end{equation}
This is the local elliptic spectral-function law
\cite{hormander-spectral,ramacher}; only its leading term at a
fixed principal point is used. To check its application, conjugation
by $\sqrt w$ preserves the full electric principal symbol and adds
only a smooth potential to the Haar Laplacian. Gauge reduction on a
principal neighborhood gives an elliptic bundle operator with scalar
leading symbol. Its local heat-kernel coefficient is proportional
to the fiber identity, giving the displayed positive constant.
Equivalently the fixed-isotypic local Weyl law applies after removing
the ineffective center. No smoothness of the entire singular quotient
or uniform estimate toward its singular orbits is asserted.

Polynomial spectral growth and Schwartz decay make
\eqref{eq:full-point-source} converge strongly for every compact
smooth test, with fixed-support seminorm bounds. It has the exact
positive, state-dependent pairing
\begin{equation}\label{eq:full-point-pairing}
B_s(f,g)=\sum_{j,n}\omega_{jn}^{-2s}|b_{jn}|^2
\overline{\widehat f(\epsilon_j\omega_{jn})}
\widehat g(\epsilon_j\omega_{jn}).
\end{equation}
Finite spectral cutoffs preserve this pairing in the strong limit.
No winding operation, phase deletion, or subtraction is used.

\begin{proposition}[Full-electric point-source exclusion]\label{prop:point-exclusion}
No real s in \eqref{eq:full-point-source} gives the complete Weil
pairing. The same holds with distinct fixed Sobolev powers in the
two orthogonal channels.
\end{proposition}
\begin{proof}
Use \eqref{eq:shrinking-test} with a nonzero real profile h.
The diagonal prime terms vanish for large R and
\eqref{eq:shrinking-diag} gives
$Q[f_R]=\norm{h''}_2^2\log R+O(1)$.
Put $p=d-2s$. If $p>0$, partial summation of
\eqref{eq:point-weyl} gives weighted cumulative mass
$dC_jR^p/p+o(R^p)$. With
$\widehat f_R(\tau)=R^{-5/2}(\tau^2+1/4)\widehat h(\tau/R)$,
Stieltjes rescaling yields
\begin{equation}\label{eq:point-scaling}
B_s[f_R]\sim d(C_0+C_1)R^{p-1}
\int_0^\infty u^{p+3}|\widehat h(u)|^2\dd u.
\end{equation}
The limiting measure is $dC_ju^{p-1}\dd u$; polynomial mass
bounds and Schwartz decay control its tails. The factor $u^4$
controls the origin, and the lower powers from $1/(4R^2)$ vanish.
For $p=0$, weighted mass is $O(\log R)$; a dyadic split at R
gives $B_s[f_R]=O(R^{-1}\log R)$. For $p<0$ the measure is finite
and $\sup|\widehat f_R|=O(R^{-1/2})$ gives $O(R^{-1})$.
Thus $p>1$ gives a positive power, $p=1$ a positive constant,
and $p<1$ decay. None gives the required logarithm. Different
powers in orthogonal channels cannot cancel positive leading terms.
\end{proof}
Fixed time rescaling changes only constants. A fixed contact
$c\ip f g_2$ cannot repair this mismatch because $\norm{f_R}_2$
is bounded. The contact in \eqref{eq:full-point-pairing} is already
fixed by the actual spectral smearing; no divergent subtraction has
been supplied with a positivity theorem. This excludes the specified
Sobolev point-source family, including genuinely distributional
members, not the full many-link electric action with every possible
seed. Other singular supports, singular-orbit preparations, and
nonclassical spectral weights require separate analysis and an
independent physical prescription.

\subsection{Many source directions on one short interval}
\begin{lemma}[Uniform local test spaces]\label{lem:many-tests}
Fix an interval I of length less than $\log2$. For all large integers
N there are N-dimensional spaces of neutral tests supported in a
fixed compact subinterval of I, parametrized by $f_{N,c}$, such that
\begin{align}
c_1\norm c&\le\norm{f_{N,c}}_2\le c_2\norm c,
&Q[f_{N,c}]&\ge c_3\log N\,\norm c^2,\label{eq:many-tests}\\
\norm{f_{N,c}}_{C^m}&\le C_mN^{m+1/2}\norm c,
&\norm{f_{N,c}}_{H^s}&\le C_sN^s\norm c\quad(s\ge0).
\end{align}
\end{lemma}
\begin{proof}
Place N disjoint translates of $h_R=R^{1/2}h(R\cdot)$ inside I,
with $R=aN$ for a sufficiently large fixed a. Let $H_c$ be their
linear combination and put $f_{N,c}=R^{-2}\T H_c$. Disjointness gives
\begin{equation}
\norm{f_{N,c}}_2^2=
\left(\norm{h''}^2+\frac{\norm{h'}^2}{2R^2}
+\frac{\norm h^2}{16R^4}\right)\norm c^2,
\qquad \norm{H_c}_2=\norm h_2\norm c.
\end{equation}
Every prime mixed term vanishes. The Fourier projection onto
$|\tau|\le\epsilon R$ satisfies
\begin{equation}\label{eq:many-low-frequency}
\norm{P_{\le\epsilon R}f_{N,c}}_2
\le(\epsilon^2+1/(4R^2))\norm{H_c}_2.
\end{equation}
Choose $\epsilon>0$ small so this is at most half the full norm,
uniformly in c for large R. The archimedean multiplier is bounded
below globally and is at least $\log R-C_\epsilon$ on the remaining
region, which carries at least three quarters of the squared norm.
This proves the Q bound, including possible cancellation between
translates. Derivative scaling and disjointness prove the integer
Sobolev and $C^m$ bounds; interpolation gives all real $s\ge0$.
\end{proof}

\subsection{Strict-half-slab distributional preparation}
Let X denote precisely the strict-half variables in
\eqref{eq:boundary-weight}. Suppose $\mathcal Af$ is a covariant
matrix-valued distribution depending on X and not on the central
spatial links U. On each fixed arithmetic support interval assume
\begin{equation}\label{eq:buffered-input}
\mathcal A:\cD^0(I)\longrightarrow H^{-s}(X)
\text{ is continuous and linear for some finite }s=s_I.
\end{equation}
This allows finite-order distributions in both the observable
variables and arithmetic input. A jointly distributional kernel
on $I\times X$ has this property by its local finite-order estimate
and Sobolev embedding in X. It includes positive-time observables
with admissible reference-color transport and finite-order
insertions, but excludes explicit central-link source dependence.

Physical boundary descent fixes
\begin{equation}\label{eq:buffered-source}
I_f(U)=\langle\mathcal Af,e^{-S_+(U,\cdot)}\rangle_X,
\qquad Jf(U)=K_\nu\mathcal Af(U)=h(U)^{-1}I_f(U).
\end{equation}
Its exact reflected pairing is
\begin{equation}\label{eq:buffered-pairing}
B_{\mathcal A}(f,g)=Z^{-1}\int_M e^{-S_0(U)}
\tfrac12\tr\big(I_f(U)^\dagger I_g(U)\big)\dd m(U).
\end{equation}
This retains the full interacting action and source dependence.
It is positive with no subtraction, and has no native-generator
or winding-module hypothesis.

The compact finite integral has a smoothing property:
\begin{equation}\label{eq:buffered-smoothing}
K_\nu:H^{-s}(X)\longrightarrow\cK_r
\text{ is bounded for every finite }s,r\ge0.
\end{equation}
Indeed every U derivative of its kernel $e^{-S_+(U,X)}/h(U)$
has uniformly bounded $H^s(X)$ norm. Duality bounds all output
derivatives, and elliptic Sobolev equivalence gives
\eqref{eq:buffered-smoothing}. Smooth gauge-compatible regularizations
of each input distribution converge in $H^{-s}(X)$ and hence strongly
in the physical norm. Thus \eqref{eq:buffered-pairing} is the actual
positive completed pairing, not a norm retained after weak escape.

\begin{theorem}[No fixed smoothing preparation]\label{thm:smoothing-exclusion}
No linear source with exact Weil pairing can obey, on one fixed
short interval I,
\begin{equation}\label{eq:common-input-order}
\norm{Jf}_{\cK_r}\le C_{I,r}\norm f_{C^m}\quad\text{for every }r>0
\end{equation}
with one finite m independent of r. In particular the physical
source \eqref{eq:buffered-input}--\eqref{eq:buffered-source} cannot
have the Weil pairing.
\end{theorem}
\begin{proof}
A conservative electric count is
\begin{equation}\label{eq:ambient-count}
\operatorname{rank}P_L\le C(1+L)^{D/2},\qquad
P_L=\mathbf1_{[0,L]}(E),\quad D=1944.
\end{equation}
The upper/lower weight bounds compare Rayleigh quotients with Haar.
On one SU(2) link the eigenvalues are $n^2-1$ with multiplicity $n^2$.
Bounding each product factor by $\sum_{n\le\sqrt{CL+1}}n^2$ and
including the four-dimensional matrix fiber proves the estimate
on the full space. Restriction to covariant functions reduces the
count. No global smooth-quotient Weyl law is needed.

Choose $L_N\asymp N^{2/D}$ so $\operatorname{rank}P_{L_N}<N$.
There is a unit coefficient vector $c_N$ in the kernel of
$c\mapsto P_{L_N}Jf_{N,c}$. Lemma~\ref{lem:many-tests} and the
electric spectral tail give
\begin{equation}\label{eq:smoothing-contradiction}
\sqrt{c_3\log N}\le\norm{Jf_{N,c_N}}
\le(1+L_N)^{-r/2}\norm{Jf_{N,c_N}}_{\cK_r}
\le C_rN^{m+1/2-r/D}.
\end{equation}
Taking $r>D(m+1/2)$ is a contradiction. For the buffered source,
continuity in \eqref{eq:buffered-input} bounds its input by one
$C^m$ seminorm. Equation~\eqref{eq:buffered-smoothing} then gives
that same input order for every r, as required.
\end{proof}
The theorem is stronger than the ordinary-input-norm exclusion and
does not require arithmetic translation to be a native action.
It also covers $J=e^{-\beta E}\mathcal A$ for fixed $\beta>0$ and
any continuous Hilbert-valued input distribution. This is electric
diffusion, not an identification with physical time; the last bound
improves to $CN^{m+1/2}\exp(-c\beta N^{2/D})$.

Any finite winding operations on X preserving
\eqref{eq:buffered-input} are included with their full observables.
Central-link windings fall outside the stated source relation; no
rational phase is removed. Fixed contact normalization cannot repair
\eqref{eq:smoothing-contradiction}, since an added fixed multiple of
$\norm f_2^2$ is bounded on the test spaces relative to $\norm c^2$.

\subsection{Necessary preparation cost and scope}
For any exact source taking values in $\cK_r$, the same rank argument
implies
\begin{equation}\label{eq:preparation-cost}
\sup_{\norm c=1}\norm{Jf_{N,c}}_{\cK_r}
\ge c_rN^{r/D}\sqrt{\log N}.
\end{equation}
Consequently a bound by the input $H^s(\R)$ norm requires
$s>r/D$, including failure at equality. The compact-occurrence
criterion of Theorem~\ref{thm:compact-occurrence} uses one fixed r
and permits a suitable input seminorm, so remains valid. Its
feasibility cannot be supplied by the fixed smoothing preparation
just excluded. Merely having smooth source vectors, with input
orders allowed to increase with output order, is not excluded.

Central-slice-dependent distributional sources outside the point
family, physically prescribed nonclassical electric weights, and
limits losing these uniform preparation estimates remain open.
The finite-lattice smoothing constants and dimension need not be
uniform in a continuum limit. Neither this theorem nor the
point-source exclusion closes all YM approaches.
